\section{Introdução e Fundamentação Teórica}\label{sec:introduction}

Bolhas de Plasma Equatoriais são grandes regiões de baixa densidade na ionosfera que ocorrem após o pôr do sol. Elas interferem gravemente nos sinais de rádio, prejudicando sistemas de navegação por satélite e comunicações \cite{zhong2025automatic, okoh2025bootstrapping}. Por isso, o problema do monitoramento contínuo desse fenômeno é fundamental para a segurança tecnológica.

A fundamentação teórica deste projeto apoia-se em estudos recentes de clima espacial e visão computacional. Tradicionalmente, a detecção do fenômeno baseia-se na análise manual de imagens do céu noturno capturadas por câmeras especiais. Contudo, esse processo exige muito tempo e esforço, tornando-se inviável diante da enorme quantidade de dados atuais \cite{githio2024machine}. Embora técnicas modernas de inteligência artificial venham sendo testadas, a literatura demonstra que o uso direto de imagens brutas exige alto poder de processamento e frequentemente confunde os algoritmos com ruídos visuais do ambiente, como nuvens e a luz da lua.

Justifica-se, portanto, a necessidade de propor o desenvolvimento de um método de inteligência artificial delimitado e específico para imagens noturnas. A proposta prepara as fotografias para remover as luzes de fundo, reduz a complexidade dos dados de forma inteligente e combina algoritmos avançados para identificar as bolhas com precisão e velocidade, viabilizando o monitoramento em tempo real.